\chapter{Measuring Market Making and Execution}\label{ch:rs:measuring-market-making-and-execution}

A market maker quoting one lot on each side of a \$100 stock ends its year \$945\,000 down. Its books show \omterm{def:rs:measuring-market-making-and-execution:pnl}{spread capture} of $+\$480\,000$, adverse selection
of $-\$1.23$ million, inventory losses of $-\$145\,000$ and fees of $-\$56\,000$. The year is twelve simulated hours of \texttt{firm.tape} scaled to 252 days
of 6.5 hours, and the decomposition is the point: it says that the quotes earn their half spread and then lose more than twice as much to the traders who
hit them, within twenty seconds; that 42\% of the fills met informed orders, which cost 1.85 ticks a share, while the uninformed fills made money; and that a
filter which trades 59\% less cuts the loss by 61\% without making a single fill much better. This chapter measures the fills of a market maker and of an
execution algorithm, with \texttt{firm.markout}: \omterm{def:rs:measuring-market-making-and-execution:curve}{mark-out curves}, fill statistics, a P\&L decomposition that adds up exactly, and \omterm{def:rs:measuring-market-making-and-execution:tca}{transaction cost analysis}.

\section{Mark-out curves}

\begin{definition}[Mark-out curve]\label{def:rs:measuring-market-making-and-execution:curve}
The \emph{mark-out curve}\index{mark-out curve} of a set of fills is the quantity-weighted average, as a function of the horizon $h$, of each fill's mark-out
(Book~2, chapter~15), $s_i\,(m(t_i + h) - p_i)$ for a fill of side $s_i$ at price $p_i$ and time $t_i$, against a reference price $m$: the mid, or the
\omterm{def:rs:order-book-features:micro}{microprice} (chapter~8).
\end{definition}

The chapter's market maker trades inside \texttt{firm.tape} as chapter~19's live stand-in: one lot on the best bid and one on the best ask, within five lots
of inventory, with latencies it draws itself (market data 0.02 seconds plus an exponential of mean 0.02, order entry 0.03 plus 0.02) and a fee of 0.05 tick a
share. Over twelve one-hour sessions it posts 18\,000 lots and 8\,144 are filled. Each fill's counterparty is known, because the tape records whether the
market order that hit the quote was informed. \Cref{fig:rs:measuring-market-making-and-execution:curve} shows the mark-outs: 0.43 tick a share at the fill
against the mid just before it (the half spread, less when the spread was one tick and the fill came at a stale level), falling to $-0.67$ at twenty seconds.
Uninformed fills keep a positive mark-out, 0.18 at twenty seconds and 0.52 at five minutes; informed ones fall to $-1.85$ and $-2.71$. Against the \omterm{def:rs:order-book-features:micro}{microprice}
the fill earns only 0.19 tick at $h = 0$: the \omterm{def:rs:order-book-features:micro}{microprice} already leans towards where the mid is going, and measures the spread that was really captured.

The curve's horizon matters as much as its level. Up to a few seconds the mark-out mixes spread and noise; beyond some horizon it stops moving, because the
information in the trade has been absorbed. For the chapter's quoter the pooled curve stays within two standard errors of its five-minute value from twenty
seconds on: the mark-outs settle at twenty seconds, and that is the horizon at which adverse selection is measured. Brogaard, Hendershott and Riordan found the
same shape in real data: high-frequency traders' liquidity-supplying orders are adversely selected, and their trading predicts price changes over horizons of
seconds.

\begin{omfigure}
\begin{tikzpicture}
\begin{semilogxaxis}[width=11.5cm,height=6cm,xlabel={horizon (seconds; the first point is the fill)},ylabel={mark-out (ticks a share)},
  tick label style={font=\small},label style={font=\small},xmin=0.04,xmax=400,ymin=-3,ymax=1,grid=major,grid style={gray!25},
  xtick={0.05,0.1,1,10,100},xticklabels={fill,0.1,1,10,100},legend columns=2,
  legend style={font=\footnotesize,at={(0.5,-0.3)},anchor=north,draw=none,fill=none,/tikz/every even column/.append style={column sep=8pt}},
  table/col sep=comma]
\addplot[name path=lo,draw=none,forget plot] table[x=h,y=all_lo]{figdata/research/23-measuring-market-making-and-execution/markout.csv};
\addplot[name path=hi,draw=none,forget plot] table[x=h,y=all_hi]{figdata/research/23-measuring-market-making-and-execution/markout.csv};
\addplot[omDef!20,forget plot] fill between[of=lo and hi];
\addplot[omDef,very thick,mark=*,mark size=1.2pt] table[x=h,y=all]{figdata/research/23-measuring-market-making-and-execution/markout.csv};
\addplot[omDef,thick,dashed] table[x=h,y=micro]{figdata/research/23-measuring-market-making-and-execution/markout.csv};
\addplot[omThm,thick,mark=triangle*,mark size=1.4pt] table[x=h,y=informed]{figdata/research/23-measuring-market-making-and-execution/markout.csv};
\addplot[black!60,thick,mark=square*,mark size=1pt] table[x=h,y=uninformed]{figdata/research/23-measuring-market-making-and-execution/markout.csv};
\draw[black!50,densely dotted] (axis cs:20,-3) -- (axis cs:20,1);
\legend{all fills (mid),all fills (microprice),informed counterparty,uninformed counterparty}
\end{semilogxaxis}
\end{tikzpicture}

\omcaption{\omterm{def:rs:measuring-market-making-and-execution:curve}{Mark-out curves} of the touch quoter's fills over twelve simulated hours, with a two-standard-error band for all fills. The dotted line marks twenty
seconds, where the curve settles. Data: \texttt{rs\_markout.curves}.}\label{fig:rs:measuring-market-making-and-execution:curve}
\end{omfigure}

\section{Fill rates and hit ratios}

\begin{definition}[Fill rate, hit ratio]\label{def:rs:measuring-market-making-and-execution:fill}
The \emph{fill rate}\index{fill rate} of a set of passive orders is the quantity filled over the quantity posted. The \emph{hit ratio}\index{hit ratio} of a
dealer answering requests for quotes is the share of its quotes that the clients trade on.
\end{definition}

A \omterm{def:rs:measuring-market-making-and-execution:fill}{fill rate} says how much of what the market maker offered was taken, and by itself cannot say whether that is good: quotes are taken most readily when they
are wrong. The chapter's second quoter withdraws its bid whenever the bid queue holds less than 30\% of the size at the touch, and its ask symmetrically,
because a thin queue predicts that the price will move through it (chapter~8):

\begin{center}
\small
\begin{tabular}{@{}lrr@{}}
\toprule
twelve simulated hours & touch quoter & with the imbalance filter\\
\midrule
lots posted, lots filled & 18\,000, 8\,144 & 16\,391, 3\,309\\
\omterm{def:rs:measuring-market-making-and-execution:fill}{fill rate} & 45\% & 20\%\\
share of fills against informed orders & 42\% & 40\%\\
mark-out at the fill, against the mid & 0.43 & 0.52\\
\quad against the \omterm{def:rs:order-book-features:micro}{microprice} & 0.19 & 0.44\\
mark-out at twenty seconds & $-0.67$ & $-0.53$\\
P\&L per lot filled & $-\$0.85$ & $-\$0.81$\\
\bottomrule
\end{tabular}
\end{center}

The filter fills less than half as often, and its fills are made when the \omterm{def:rs:order-book-features:micro}{microprice} sits near the mid, so it captures more of the spread (0.44 tick against
the \omterm{def:rs:order-book-features:micro}{microprice} rather than 0.19); but the traders who hit it are as often informed as before, and per lot it loses nearly as much. A dealer's \omterm{def:rs:measuring-market-making-and-execution:fill}{hit ratio} has
the same ambiguity on a request-for-quote platform: a dealer who wins every request is quoting too well.

\section{The market maker's P\&L, decomposed}

\begin{definition}[Spread capture, adverse-selection cost, inventory P\&L]\label{def:rs:measuring-market-making-and-execution:pnl}
For a fill of side $s$, quantity $q$ and price $p$ at time $t$, with reference price $m$, horizon $H$ and the end of the period $T$: the \emph{spread
capture}\index{spread capture} is $sq\,(m(t) - p)$; the \emph{adverse-selection cost}\index{adverse-selection cost} is $sq\,(m(t + H) - m(t))$; the \emph{inventory
P\&L}\index{inventory P\&L} is $sq\,(m(T) - m(t + H))$. Summed over fills, with fees, the three add up to the P\&L of the fills with the final position marked at
$m(T)$.
\end{definition}

The identity is exact: $m(T) - p = (m(t) - p) + (m(t+H) - m(t)) + (m(T) - m(t+H))$ for each fill (\cref{lst:rs:measuring-market-making-and-execution:decompose}).
What it leaves to judgement is the horizon $H$: too short, and the adverse selection still unfolding after it is booked as inventory; too long, and inventory
risk the market maker chose to carry is booked as adverse selection. The settling horizon of the \omterm{def:rs:measuring-market-making-and-execution:curve}{mark-out curve} is the natural choice, twenty seconds here.
Menkveld's decomposition of a real high-frequency market maker on Chi-X and Euronext has the same structure: a profit of \texteuro1.55 a trade on the spread net
of fees, a positioning loss of \texteuro0.68 (a profit on positions held less than five seconds, a loss on longer ones), a gross profit of \texteuro0.88.

\begin{omfigure}
\begin{tikzpicture}
\begin{axis}[width=11cm,height=5.8cm,ybar,bar width=12pt,ylabel={dollars, twelve hours},symbolic x coords={spread,adverse,inventory,fees,total},
  xtick=data,xticklabels={spread\\capture,adverse\\selection,inventory,fees,total},xticklabel style={align=center},
  tick label style={font=\small},label style={font=\small},enlarge x limits=0.15,
  ymin=-10000,ymax=5000,grid=major,grid style={gray!25},legend columns=2,
  yticklabel style={/pgf/number format/fixed,/pgf/number format/1000 sep={\,}},scaled y ticks=false,
  legend style={font=\footnotesize,at={(0.5,-0.3)},anchor=north,draw=none,fill=none,/tikz/every even column/.append style={column sep=8pt}},
  table/col sep=comma]
\addplot[fill=omDef!55,draw=omDef] table[x=part,y=plain]{figdata/research/23-measuring-market-making-and-execution/decomp.csv};
\addplot[fill=gray!45,draw=gray] table[x=part,y=pull]{figdata/research/23-measuring-market-making-and-execution/decomp.csv};
\legend{touch quoter,with the imbalance filter}
\end{axis}
\end{tikzpicture}

\omcaption{The P\&L of the two quoters over twelve simulated hours, decomposed with a horizon of twenty seconds. Data:
\texttt{rs\_markout.decomposition}.}\label{fig:rs:measuring-market-making-and-execution:decomp}
\end{omfigure}

Over the twelve hours the touch quoter captures \$3\,519 of spread, loses \$8\,977 to adverse selection within twenty seconds, \$1\,060 on inventory after it, and
pays \$407 of fees: a loss of \$6\,925 (\cref{fig:rs:measuring-market-making-and-execution:decomp}). Per share, the spread brings 0.43 tick and adverse selection
takes 1.10, which is the mark-out of $-0.67$ at twenty seconds. The filtered quoter captures \$1\,734, loses \$3\,483 and \$771, pays \$165: $-\$2\,685$. Neither
is a business; the decomposition says what would have to change (fewer informed counterparties, or a wider spread) and that inventory is not the problem.

\section{Measuring execution}

\begin{definition}[Delay cost, opportunity cost, VWAP slippage]\label{def:rs:measuring-market-making-and-execution:costs}
For a parent order decided at price $P_d$, released at the \omterm{def:rs:simulation-versus-live:calibration}{arrival price} $P_a$, filled for $Q$ of its target $X$ at average price $\bar P$, with the price $P_e$
at its end: the \emph{delay cost}\index{delay cost} is $sQ(P_a - P_d)$, the execution cost $sQ(\bar P - P_a)$, the \emph{opportunity cost}\index{opportunity cost}
$s(X - Q)(P_e - P_d)$; with fees they add up to the \omterm{def:rs:simulation-versus-live:calibration}{implementation shortfall} (chapter~19). The \emph{VWAP slippage}\index{VWAP slippage} is $s(\bar P - \mathrm{VWAP})$,
the average price against the market's \omterm{def:rs:market-data-for-research:bar}{volume-weighted average price} over the order's life.
\end{definition}

The chapter's parent order buys 300 lots: decided at minute 10, released at minute 11, due by minute 30. Every twenty seconds the executor rests the next step's
lots at the best bid and sends a market order for whatever is more than eight lots behind a straight-line schedule. In one session (seed 61) it fills 28\,700
shares, 43\% passively on average over the sessions, and its report reads:

\begin{center}
\small
\begin{tabular}{@{}lr@{}}
\toprule
\omterm{def:rs:measuring-market-making-and-execution:tca}{transaction cost analysis}, one order & ticks a share\\
\midrule
\omterm{def:rs:measuring-market-making-and-execution:costs}{delay cost} & 0.00\\
execution cost against arrival & $-7.83$\\
\omterm{def:rs:measuring-market-making-and-execution:costs}{opportunity cost} (1\,300 shares unfilled) & $-0.91$\\
fees & 0.05\\
\omterm{def:rs:simulation-versus-live:calibration}{implementation shortfall} & $-8.69$\\
\omterm{def:rs:measuring-market-making-and-execution:costs}{VWAP slippage} & $-1.10$\\
\bottomrule
\end{tabular}
\end{center}

A negative shortfall of 8.69 ticks (8.7 basis points at \$100) looks like a triumph of execution. It is the market: the mid fell over those nineteen minutes, and
a buyer who arrives before a fall looks brilliant against arrival. Over the twelve sessions the shortfall averages $-2.00$ ticks a share with a standard error of
2.73; its execution part, $-4.68$ (2.09), splits into the drift of the mid from arrival to each fill, $-4.78$ (2.07), and the price paid against the mid just
before the fill, $+0.10$ (0.03). Only the second is the algorithm's doing, and only it is measured precisely. Even the \omterm{def:rs:measuring-market-making-and-execution:costs}{delay cost}, $+2.83$ (0.94), is chance: the
same twelve sessions run without the order rise by 2.92 ticks in that minute. And the simulator allows what no desk can do: run the same session without the
order. The mid's move from minute 11 to minute 30 is $-9.71$ ticks with the order and $-9.96$ without it, an impact of $+0.25$ (0.24). In this market the efficient
price does not respond to the order, so its lasting impact is nil by construction; a real market's would not be (chapter~27).

\begin{omfigure}
\begin{tikzpicture}
\begin{axis}[width=9.5cm,height=6cm,xlabel={mid move, minute 11 to 30, without the order (ticks)},ylabel={ticks a share},tick label style={font=\small},
  label style={font=\small},xmin=-60,xmax=20,ymin=-25,ymax=12,grid=major,grid style={gray!25},legend pos=north west,legend style={font=\footnotesize},
  table/col sep=comma]
\addplot[only marks,mark=*,mark size=2pt,omDef] table[x=move_without,y=execution]{figdata/research/23-measuring-market-making-and-execution/tca.csv};
\addplot[only marks,mark=square*,mark size=1.8pt,omThm] table[x=move_without,y=paid]{figdata/research/23-measuring-market-making-and-execution/tca.csv};
\legend{execution cost against arrival,price paid against the mid}
\end{axis}
\end{tikzpicture}

\omcaption{Twelve parent orders: the execution cost measured against the \omterm{def:rs:simulation-versus-live:calibration}{arrival price} follows the market's own move over the order's life (measured in the same
session run without the order); the price paid against the mid at each fill does not. Data: \texttt{rs\_markout.tca\_all}.}\label{fig:rs:measuring-market-making-and-execution:tca}
\end{omfigure}

\section{Transaction cost analysis}

\begin{definition}[Transaction cost analysis]\label{def:rs:measuring-market-making-and-execution:tca}
\emph{Transaction cost analysis}\index{transaction cost analysis} (TCA) is the measurement of the costs of executed orders against benchmarks (the decision and
\omterm{def:rs:simulation-versus-live:calibration}{arrival prices}, VWAP, the close), decomposed into parts attributable to decisions (delay, the choice to leave quantity unfilled) and to execution (the price paid
against the market at each fill), and aggregated over many orders to compare brokers, algorithms and traders.
\end{definition}

The lesson of the twelve orders is that TCA is a statistical exercise with a large noise term. An arrival-price shortfall has the variance of the market's move
over the order's life, which dwarfs the cost being measured; averages over many orders, conditioning on the market's move (regressing shortfall on the
contemporaneous return of the stock or its sector), and cost measures anchored at each fill (the price paid against the mid, the fill's own mark-out) are what
make it informative. \omterm{def:rs:measuring-market-making-and-execution:costs}{VWAP slippage} has less noise (0.40 tick with a standard error of 0.61 here) because the benchmark moves with the market, and it can be gamed
for the same reason: an algorithm that trades with the volume matches VWAP whatever it costs. A TCA report should state the benchmark, the decomposition, the
standard error, and the number of orders, and compare algorithms by randomised experiment (chapter~21).

\section{Tutorial: the flat year}\label{tut:rs:measuring-market-making-and-execution:tut}

\begin{tutorial}
\textbf{Goal.} Run two market makers and an executor inside \texttt{firm.tape}, draw the \omterm{def:rs:measuring-market-making-and-execution:curve}{mark-out curves} by counterparty, decompose the P\&L, and write the
TCA of a parent order. \textbf{End state:} the two tables and
\cref{fig:rs:measuring-market-making-and-execution:curve,fig:rs:measuring-market-making-and-execution:decomp,fig:rs:measuring-market-making-and-execution:tca}.
\begin{tutsteps}
\item \textbf{Mark-outs}: the reference in force at each time, and every fill against it at every horizon.
\omcode{code/firm/markout/firm_markout.py}{38}{45}{Mark-outs against a reference price path.}\label{lst:rs:measuring-market-making-and-execution:markouts}
\item \textbf{The decomposition} that adds up.
\omcode{code/firm/markout/firm_markout.py}{77}{90}{Spread capture, adverse selection, inventory, fees.}\label{lst:rs:measuring-market-making-and-execution:decompose}
\item \textbf{The shortfall} of a parent order.
\omcode{code/firm/markout/firm_markout.py}{93}{103}{Implementation shortfall in four parts.}\label{lst:rs:measuring-market-making-and-execution:shortfall}
\item \textbf{Run} \texttt{rs\_markout.curves}, \texttt{decomposition}, \texttt{settle}, \texttt{execution}, \texttt{tca\_all} and \texttt{fig\_markout.py}.
\end{tutsteps}
\textbf{What to change next.} Quote one tick behind the touch and see the \omterm{def:rs:measuring-market-making-and-execution:fill}{fill rate} fall and the per-lot mark-out rise; give the executor a signal (the
\omterm{def:rs:order-book-features:micro}{microprice}) and measure whether its price paid against the mid improves.
\end{tutorial}

\section{Build: the measurement module}\label{bld:rs:measuring-market-making-and-execution:markout}

\begin{build}
\bfield{Purpose} Every fill the firm makes, passive or aggressive, is measured the same way: mark-outs at standard horizons against the mid and the \omterm{def:rs:order-book-features:micro}{microprice},
grouped by counterparty, venue and strategy; market-making P\&L decomposed; every parent order analysed.
\bfield{Interface} \texttt{microprice}, \texttt{ref\_at(ref\_t, ref\_px, t)}, \texttt{markouts(t, side, px, ref\_t, ref\_px, horizons)}, \texttt{curve(M, qty,
groups)}, \texttt{settle\_horizon(mean, horizons, tol)}, \texttt{fill\_rate}, \texttt{hit\_ratio}, \texttt{mm\_decompose(t, side, px, qty, ref\_t, ref\_px, H, fee,
t\_end)}, \texttt{shortfall(side, target, decision\_px, arrival\_px, fills, end\_px, fee)}, \texttt{vwap\_slippage}, \texttt{tca\_report}.
\bfield{Rules} The reference is the one in force before the fill; the decomposition adds up to the marked P\&L to the cent; horizons are fixed firm-wide;
every average is reported with its standard error and its count.
\bfield{Acceptance tests} \texttt{code/firm/markout/tests/}: references and mark-outs by hand; a weighted curve; the settling horizon; the decomposition equal to
the P\&L of random fills with the position marked at the end; the shortfall equal to the paper portfolio's gain minus the real one's; \omterm{def:rs:measuring-market-making-and-execution:costs}{VWAP slippage}, \omterm{def:rs:measuring-market-making-and-execution:fill}{fill rate} and
\omterm{def:rs:measuring-market-making-and-execution:fill}{hit ratio} by hand.
\bfield{Stretch} Mark-outs against a fair-value model; regression-adjusted TCA (shortfall on the market's move); counterparty scorecards for a dealer.
\end{build}

\begin{omsources}
\item A.~J.~Menkveld, ``High frequency trading and the new market makers'', \emph{Journal of Financial Markets} 16(4), 2013 (Tinbergen Institute discussion
paper 11-076, 2011).
\item J.~Brogaard, T.~Hendershott and R.~Riordan, ``High-frequency trading and price discovery'', \emph{Review of Financial Studies} 27(8), 2014.
\item A.~F.~Perold, ``The implementation shortfall: paper versus reality'', \emph{Journal of Portfolio Management} 14(3), 1988.
\end{omsources}

\section{Exercises}

\begin{exercise}[$\star$]\label{exo:rs:measuring-market-making-and-execution:1}
A buy fill at 99.99 when the mid was 100.00; twenty seconds later the mid is 99.97, at the end of the day 100.02. Split the fill's P\&L per share into \omterm{def:rs:measuring-market-making-and-execution:pnl}{spread
capture}, adverse selection and inventory (in ticks of 0.01).
\end{exercise}

\begin{exercise}[$\star$]\label{exo:rs:measuring-market-making-and-execution:2}
The touch quoter posted 18\,000 lots and 8\,144 were filled. What is its \omterm{def:rs:measuring-market-making-and-execution:fill}{fill rate}? Why is a higher \omterm{def:rs:measuring-market-making-and-execution:fill}{fill rate} not better?
\end{exercise}

\begin{exercise}[$\star$]\label{exo:rs:measuring-market-making-and-execution:3}
The spread brings 0.43 tick a share and adverse selection to twenty seconds takes 1.10. What is the mark-out at twenty seconds?
\end{exercise}

\begin{exercise}[$\star\star$]\label{exo:rs:measuring-market-making-and-execution:4}
Why does the choice of the horizon $H$ not change the total of the decomposition? Which components does it move, and in which direction when $H$ is too short?
\end{exercise}

\begin{exercise}[$\star\star$]\label{exo:rs:measuring-market-making-and-execution:5}
A buy order of 100 shares is decided at 50.00, released at 50.20, filled 80 shares at an average of 50.40; the price is 51.00 at its end; fees are 0.01 a share.
Compute the delay, execution and \omterm{def:rs:measuring-market-making-and-execution:costs}{opportunity costs} and the shortfall.
\end{exercise}

\begin{exercise}[$\star\star$]\label{exo:rs:measuring-market-making-and-execution:6}
The imbalance filter cut the loss from \$6\,925 to \$2\,685. How much of that is fewer fills and how much better fills?
\end{exercise}

\begin{exercise}[$\star\star\star$]\label{exo:rs:measuring-market-making-and-execution:7}
\emph{Coding.} Scale the chapter's twelve-hour decomposition of the touch quoter to a year of 252 days of 6.5 hours, and state what the scaling assumes.
\end{exercise}

\begin{exercise}[$\star\star\star$]\label{exo:rs:measuring-market-making-and-execution:8}
\emph{Find the flaw.} ``Our new execution algorithm beat arrival by 8.7 basis points on its first order, so it is the best algorithm we have.''
\end{exercise}

\section{Problem: The Flat Year}

\begin{problem}[{Weekend problem --- the market maker's books}]\label{pb:rs:measuring-market-making-and-execution:1}
The chapter's two quoters and its executor, in twelve simulated hours of \texttt{firm.tape}.

\medskip\noindent\textbf{Part I --- Mark-outs.}
\begin{enumerate}
\item What is the touch quoter's mark-out at the fill, against the mid and against the \omterm{def:rs:order-book-features:micro}{microprice}, and why do they differ?
\item What are the informed and uninformed mark-outs at twenty seconds and five minutes?
\item What share of the fills met informed orders?
\item At which horizon does the curve settle, by which rule?
\end{enumerate}

\medskip\noindent\textbf{Part II --- Fills.}
\begin{enumerate}[resume]
\item Give both quoters' lots posted, lots filled and \omterm{def:rs:measuring-market-making-and-execution:fill}{fill rates}.
\item What does the imbalance filter change in the mark-outs, and what does it not change?
\item What is each quoter's P\&L per lot?
\end{enumerate}

\medskip\noindent\textbf{Part III --- The decomposition.}
\begin{enumerate}[resume]
\item Decompose the touch quoter's twelve hours at twenty seconds.
\item Decompose the filtered quoter's.
\item How do the spread and adverse-selection components per share relate to the \omterm{def:rs:measuring-market-making-and-execution:curve}{mark-out curve}?
\item What does Menkveld's decomposition of a real market maker look like?
\end{enumerate}

\medskip\noindent\textbf{Part IV --- Execution and the verdict.}
\begin{enumerate}[resume]
\item Give the one-order TCA report of seed 61.
\item Average the shortfall and its parts over the twelve orders, with standard errors.
\item Which part is the algorithm's doing?
\item What is the order's impact on the mid, measured by the counterfactual, and why is it small here?
\item State the \emph{named result}: the touch quoter's year decomposed into \omterm{def:rs:measuring-market-making-and-execution:pnl}{spread capture}, adverse selection, inventory and fees, and the horizon at which its
mark-outs settle.
\item What would the market maker have to change to break even?
\item How should a desk compare two execution algorithms?
\item What should every TCA report state?
\item In one sentence: what does a \omterm{def:rs:measuring-market-making-and-execution:curve}{mark-out curve} measure?
\end{enumerate}
\end{problem}

\section{Interview questions}

\begin{interviewq}[$\star$ \iqroles{trader}]\label{iq:rs:measuring-market-making-and-execution:1}
What is a mark-out, and what horizon would you use?
\end{interviewq}

\begin{interviewq}[$\star\star$ \iqroles{trader, researcher}]\label{iq:rs:measuring-market-making-and-execution:2}
Decompose a market maker's P\&L. How do you choose the split between adverse selection and inventory?
\end{interviewq}

\begin{interviewq}[$\star\star$ \iqroles{researcher}]\label{iq:rs:measuring-market-making-and-execution:3}
A high-frequency market maker earns \texteuro1.55 a trade on the spread and loses \texteuro0.68 on positions. What does this tell you about its business?
\end{interviewq}

\begin{interviewq}[$\star\star$ \iqroles{trader, bank}]\label{iq:rs:measuring-market-making-and-execution:4}
Your \omterm{def:rs:measuring-market-making-and-execution:fill}{hit ratio} on a request-for-quote platform doubled last month. Good news?
\end{interviewq}

\begin{interviewq}[$\star\star$ \iqroles{trader, researcher}]\label{iq:rs:measuring-market-making-and-execution:5}
How do you tell whether an execution algorithm is better than another, given that each order's shortfall is dominated by the market's move?
\end{interviewq}

\begin{interviewq}[$\star\star\star$ \iqroles{researcher, developer}]\label{iq:rs:measuring-market-making-and-execution:6}
Design the firm's fill-measurement service: inputs, reference prices, horizons, groupings, and the checks that its numbers add up.
\end{interviewq}
